% Supplemental controlled branch study; not a full exploration result.
\subsection{Category interventions on native candidate branches}
\label{supp:native-semantic-mechanism}

\paragraph{Matched acquired state.}
A controlled branching scene tests whether category input can change the original candidate selection when visible geometry is ambiguous. Two procedural arrangements exchange an open structure and a closed structure behind partial screens. For each arrangement, one acquired cold-start state supplies the original four candidates, the shared geometric bank and the common coverage/measured-surface terms. All four native candidates lie on the right side of the scene, which limits the available observation alternatives. Correct, neutral and swapped category priors are evaluated at both $T=1$ and the independently selected $T=16$, giving 12 scoring conditions. No goal, route, common score weight or stopping rule is inserted to obtain a preferred result. The first fixture failed to produce modeled objects under the frozen frontend and remains a failed representability pilot. The revised physical fixture provides visible posts and top support; its geometry is frozen before candidate outcomes. Rendering checks retain cross-layout float32 depth differences of at most $5\times10^{-7}$\,m and corresponding template-gain differences of approximately $0.034$\,m$^2$. The two arrangements therefore do not constitute an exact sensor counterfactual. Within each arrangement, all cue and temperature conditions share exactly the same acquired input and point basis.

\paragraph{Measured consequences of selection.}
All eight original outbound routes are acquired once from their respective common initial states. They complete 114 actions in total, reach all eight endpoints and record no simulated collisions. Each branch fuses only its acquired depth into a native TSDF. Independent evaluation uses fixed source-surface references and reports both object and global surface increments, correct-free coverage and execution cost. Scoring conditions reuse these eight outcomes; they are neither 12 independent trials nor complete autonomous exploration episodes.

At $T=16$, the correct cue changes the right-arrangement selection. Object net recovered area rises from $0.046$ to $1.403$\,m$^2$, outbound actions fall from 17 to 15, and correct-free coverage gain increases by 3.33 percentage points. Global net recovered surface nevertheless decreases by $24.122$\,m$^2$. The left arrangement retains its selected route. At $T=1$, the correct left cue changes the target but reduces object recovery while increasing global surface recovery. Neutral and swapped cues, including adverse outcomes, remain in Table~\ref{tab:b9-mechanism}. These contrasts demonstrate a category-dependent native target and a measurable allocation tradeoff. They establish neither a uniform reconstruction improvement nor an unknown-exploration efficiency advantage.

\begin{table}[tbp]
\centering\small
\caption{All 12 controlled mechanism conditions. Each layout retains its four original candidate routes; eight actual outbound branches (114 actions) supply all outcomes. Columns report S minus G after selecting from the same candidate set. Surface increments use native TSDF distance at 5\,cm against fixed visible source triangles; coverage is correct-free area in percentage points. Within each layout, cue and temperature interventions share exactly the same acquired bank. The two layouts have disclosed float32 rendering differences and are not exact sensor counterfactuals. These are short fixed-route branches, not complete exploration trials.}
\label{tab:b9-mechanism}
\begin{tabular}{lllrrrrr}
\toprule
Layout & $T$ & Cue & G$\to$S & $\Delta A_{obj}$ (m$^2$) & $\Delta A_{all}$ (m$^2$) & $\Delta C$ (pp) & $\Delta$ actions \\
\midrule
Left & 1 & Correct & 0 $\to$ 2 & -1.221 & +24.534 & -2.33 & +2 \\
Left & 1 & Neutral & 0 $\to$ 0 & +0.000 & +0.000 & +0.00 & +0 \\
Left & 1 & Swapped & 0 $\to$ 0 & +0.000 & +0.000 & +0.00 & +0 \\
Left & 16 & Correct & 2 $\to$ 2 & +0.000 & +0.000 & +0.00 & +0 \\
Left & 16 & Neutral & 2 $\to$ 2 & +0.000 & +0.000 & +0.00 & +0 \\
Left & 16 & Swapped & 2 $\to$ 0 & +1.221 & -24.534 & +2.33 & -2 \\
Right & 1 & Correct & 0 $\to$ 0 & +0.000 & +0.000 & +0.00 & +0 \\
Right & 1 & Neutral & 0 $\to$ 0 & +0.000 & +0.000 & +0.00 & +0 \\
Right & 1 & Swapped & 0 $\to$ 2 & -1.356 & +24.122 & -3.33 & +2 \\
Right & 16 & Correct & 2 $\to$ 0 & +1.356 & -24.122 & +3.33 & -2 \\
Right & 16 & Neutral & 2 $\to$ 2 & +0.000 & +0.000 & +0.00 & +0 \\
Right & 16 & Swapped & 2 $\to$ 2 & +0.000 & +0.000 & +0.00 & +0 \\
\bottomrule
\end{tabular}
\end{table}

